\section{Mathematical Analysis of the Estimator}\label{sec:analysis}

\subsection{Error decomposition}

Let $u \in H^1_g(\Omega)$ be the exact solution of \eqref{eq:model-problem},
and let $\hat{u} \in H^1(\Omega)$ be a PINN approximation.
Define the error
\[
e := u - \hat{u}.
\]

Since boundary conditions are enforced weakly, $e \notin H^1_0(\Omega)$ in general.
Following Section~\ref{sec:methodology}, we introduce the decomposition
\begin{equation}\label{eq:analysis-decomposition}
e = e_0 + w,
\end{equation}
where $e_0 \in H^1_0(\Omega)$ and $w \in H^1(\Omega)$ solves
\begin{equation}\label{eq:analysis-lifting}
\begin{aligned}
- \Delta w &= 0 && \text{in } \Omega,\\
w &= g - \hat{u} && \text{on } \partial\Omega.
\end{aligned}
\end{equation}

\subsection{Residual functional}

Define the residual functional $R \in H^{-1}(\Omega)$ by
\begin{equation}\label{eq:analysis-residual}
\langle R, v \rangle
=
\int_\Omega f v\,dx
-
\int_\Omega a \nabla \hat{u} \cdot \nabla v\,dx,
\quad \forall v \in H^1_0(\Omega).
\end{equation}

This corresponds to the strong residual
\[
R = f + \nabla \cdot (a \nabla \hat{u})
\]
in the distributional sense.

\subsection{Error equation}

Subtracting the weak formulation \eqref{eq:weak-form} from
\eqref{eq:analysis-residual}, we obtain for all $v \in H^1_0(\Omega)$:
\begin{equation}\label{eq:error-equation}
\int_\Omega a \nabla e \cdot \nabla v \,dx
=
\langle R, v \rangle.
\end{equation}

Since $e = e_0 + w$ and $v \in H^1_0(\Omega)$, we have
\[
\int_\Omega a \nabla e_0 \cdot \nabla v \,dx
=
\langle R, v \rangle - \int_\Omega a \nabla w \cdot \nabla v \,dx.
\]

However, choosing $v = e_0$ yields
\begin{equation}\label{eq:error-test}
\int_\Omega a \nabla e_0 \cdot \nabla e_0 \,dx
=
\langle R, e_0 \rangle - \int_\Omega a \nabla w \cdot \nabla e_0 \,dx.
\end{equation}

To avoid coupling terms, we estimate $e_0$ directly using the residual.

\subsection{Main reliability result (constant coefficient case)}

\begin{theorem}[A posteriori reliability]\label{thm:main}
Assume $a \equiv \text{const}$. Then the PINN approximation satisfies
\begin{equation}\label{eq:main-bound}
\|\nabla(u - \hat{u})\|_{L^2(\Omega)}
\le
\frac{1}{\alpha}\|R\|_{H^{-1}(\Omega)}
+
\|\nabla w\|_{L^2(\Omega)}.
\end{equation}
\end{theorem}

\begin{proof}

Taking $v = e_0 \in H^1_0(\Omega)$ in \eqref{eq:error-equation}, we obtain
\[
\int_\Omega a \nabla e \cdot \nabla e_0
=
\langle R, e_0 \rangle.
\]

Since $e = e_0 + w$,
\[
\int_\Omega a \nabla e_0 \cdot \nabla e_0
+
\int_\Omega a \nabla w \cdot \nabla e_0
=
\langle R, e_0 \rangle.
\]

For constant $a$, orthogonality of harmonic lifting implies
\[
\int_\Omega a \nabla w \cdot \nabla e_0 = 0.
\]

Hence,
\[
\int_\Omega a \nabla e_0 \cdot \nabla e_0
=
\langle R, e_0 \rangle.
\]

Using coercivity and dual norm definition,
\[
\alpha \|\nabla e_0\|^2
\le
\|R\|_{H^{-1}} \|\nabla e_0\|.
\]

Therefore,
\[
\|\nabla e_0\|
\le
\frac{1}{\alpha} \|R\|_{H^{-1}}.
\]

Finally, by triangle inequality,
\[
\|\nabla e\|
\le
\|\nabla e_0\|
+
\|\nabla w\|,
\]
which yields the result.
\end{proof}

\subsection{Variable coefficient case}

For variable diffusion coefficients $a(\bm{x})$, the orthogonality property
used above no longer holds. The exact decomposition would require an
$a$-harmonic lifting $w_a$ satisfying
\[
- \nabla \cdot (a \nabla w_a) = 0,
\]
which is not directly available in a computationally efficient form.

Instead, we use the Laplace-based lifting $w$ defined in
\eqref{eq:analysis-lifting} as a computable surrogate.
Although this introduces a consistency gap, the resulting estimator remains
a valid and practical upper bound in all numerical experiments.

\subsection{Final estimator}

Combining the above results, we define the estimator
\begin{equation}\label{eq:estimator}
\eta :=
\frac{1}{\alpha}\|R\|_{H^{-1}(\Omega)}
+
\|\nabla w\|_{L^2(\Omega)}.
\end{equation}

\subsection{Effectivity index}

We assess the sharpness of the estimator using the effectivity index
\begin{equation}
\eta_{\mathrm{eff}} :=
\frac{\eta}{\|\nabla(u - \hat{u})\|_{L^2(\Omega)}}.
\end{equation}

The numerical results reported in Table~\ref{tab:benchmark_results}
show that
\[
1 \le \eta_{\mathrm{eff}} \le 1.31,
\]
confirming both reliability and near-optimal sharpness of the estimator
across all benchmarks. :contentReference[oaicite:0]{index=0}